\documentclass[11pt,letterpaper]{article}
\usepackage{technical-testing}
% Shared non-semantic helper definitions; one semantic module is loaded by each root.
\newcommand{\PTSource}[1]{\par\smallskip{\footnotesize\color{Meta}\textbf{Book crosswalk:} #1}\par\smallskip}
\newcommand{\PTOutcome}[1]{\begin{tipbox}[Competency outcome]#1\end{tipbox}}
\newcommand{\PTEvidence}[1]{\begin{notebox}[Required evidence]#1\end{notebox}}
\newcount\ptlinecount
\newcommand{\PTLines}[1]{\par\ptlinecount=0\loop\ifnum\ptlinecount<#1\par\noindent\rule{\linewidth}{0.25pt}\vspace{0.22cm}\advance\ptlinecount by1\repeat\par}
\newcommand{\PTField}[1]{\noindent\textbf{#1}\quad\rule{0.66\linewidth}{0.3pt}\par\vspace{0.22cm}}
\newcommand{\PTRating}{\noindent\textbf{Rating:}\quad 0\quad 1\quad 2\quad 3\quad 4\qquad\textbf{Evidence ref:}\rule{0.30\linewidth}{0.3pt}\par\vspace{0.14cm}}
\newcommand{\PTTask}[2]{\par\medskip\noindent\textbf{Task #1.} #2\par}

\setDocTitle{Parser Verification and Release Review Procedure}
\setDocSubtitle{Correctness, language coverage, safety and reproducibility gates}
\setDocVersion{1.0}\setDocStatus{Engineering education / source-aligned}
\setDocOwner{Compiler Engineering}
\setDocDate{30 October 2026}
\setDocSummary{Engineering verification procedure.}
\begin{document}\makedoctitle\makedocsummary
\tableofcontents\clearpage
\section{Review sequence}
\begin{enumerate}\item Validate tokens and source encoding; record complete/partial-input policy.\item Inspect grammar for undefined, unreachable, nonproductive, nullable and cyclic forms.\item Establish derivation acceptance and structural output independently.\item Run automated positive, negative and metamorphic tests.\item Validate lexer/parser error spans and distinct diagnostic categories.\item Enforce CPU/space, nesting, token-count and input-size controls.\item For C++, run ASan/UBSan where toolchain supports them.\item Reproduce performance result with raw data and revision pin.\item Record scorecard review, unresolved defects and release decision.\end{enumerate}
\section{Traceability}
Maintain stable IDs across grammar rules, threat scenarios, tests, and defects; link to the original input and parser revision. Negative results are evidence only for the tested environment, not proof of absence of vulnerabilities.
\section{Failure handling}
Block release for infinite loops, crashes on malformed input, silent semantic execution after recovery, and unbounded resource exhaustion within the documented threat model. Remediation requires a minimized regression and a passing repeat run.

\end{document}
